\subsection{Implementations}

\subsubsection{AddressSanitizer and Miri}

AddressSanitizer~\cite{serebryany2012addresssanitizer} is a memory bug detector
based on instrumentation, and
%
Miri~\cite{jung2026miri, gh-miri} is a Rust interpreter designed to detect
undefined behavior, including memory bugs.
%
They can serve as test oracles by checking whether a program that passes type
checking incurs a memory bug during execution.
%
However, since type checking is outside their scope, they cannot identify
soundness bugs that do not manifest as undefined behavior.

To empirically validate their effectiveness, we applied them to the code from
the eleven soundness issues that result in memory bugs.
%
In our experiments, both AddressSanitizer and Miri successfully detected memory
bugs in all cases.
%
Therefore, using AddressSanitizer or Miri as a test oracle is a promising
direction, although complementary tools are necessary to cover a broader range
of soundness bugs.

\subsubsection{Chalk}

Chalk~\cite{gh-chalk} is an experimental trait solver that serves as a
prototyping ground for designing and improving the trait solver in \rustc{}.
%
It determines whether certain \C{impl}s are overlapping or orphan and whether a
given type implements a specific trait.
%
Chalk approaches these tasks in a fundamentally different way from \rustc{};
%
it models Rust declarations using logical predicates and performs checks by
solving logical formulas.
%
Thus, it is reasonable to utilize Chalk as a test oracle to discover
trait-related soundness bugs.
%
However, since Chalk targets only the trait system, other type system features
remain outside its scope.

To assess its effectiveness, we applied the latest version (0.104.0) to the
seven trait-affecting issues (99554, 105787, 106040, 114061, 114389, 126079, and
133361).
%
% As Chalk's syntax differs from Rust's, we manually rewrote each code snippet
% in Chalk's syntax.
%
Unfortunately, Chalk currently cannot express Issue 106040 because it does not
support associated types with bounds in \C{impl}s.
%
For the remaining issues, Chalk correctly identified the orphan \C{impl} in
Issue 99554 and the dyn-compatibility violation in Issue 126079;
%
it also rejected Issue 105787, but the message incorrectly cited an orphan
\C{impl} rather than overlapping \C{impl}s;
%
it crashed on Issue 114061 and accepted Issues 114389 and 133361.
%
Consequently, we conclude that Chalk is yet immature for use as a practical test
oracle.

To facilitate its realistic use as an oracle for trait-related soundness bugs,
we suggest improving its robustness and correctness by testing it with diverse
inputs.
%
As Chalk targets only the trait system, we do not consider it a general oracle
that can cover all kinds of soundness bugs, and we do not suggest extending it
to cover features irrelevant to traits.

\subsubsection{a-mir-formality}

\begin{figure}[t]
% \begin{lstlisting}[caption={Issue 99554 in Chalk}, label={lst:chalk}]
% #[upstream] trait X<T, U> {}
% struct S {}
% trait Y { type A; }
% impl<T> Y for T { type A = T; }
% impl<T> X<S, T> for <T as Y>::A where T: Y {}
% \end{lstlisting}
% \vspace{-1em}
\begin{lstlisting}[caption={Issue 99554 in a-mir-formality}, label={lst:a-mir-formality}]
[ crate a { trait X<ty T, ty U> {} },
  crate b { struct S {}
            trait Y { type A: []; }
            impl<ty T> Y for T { type A = T; }
            impl<ty T> X<S, T> for <T as Y>::A where T: Y {} } ]
\end{lstlisting}
\vspace{-1em}
\end{figure}

The a-mir-formality project~\cite{gh-amirformality} aims to develop an
executable formal model of the Rust type system, which can decide whether a
given program should pass type checking or not.
%
As its scope is not limited to a specific part of the type system, we believe it
has strong potential to become a comprehensive test oracle for type soundness.

To validate its effectiveness, we applied a-mir-formality at commit \C{ab4f18e}
to the collected soundness issues.
%
As a-mir-formality defines its own syntax, we manually rewrote the code;
%
for top-level items, its syntax is close to Rust's;
%
for function bodies, it follows MiniRust~\cite{gh-minirust}, which resembles the
mid-level intermediate representation (MIR) used in \rustc{}.
%
For example, \cref{lst:a-mir-formality} shows the code for Issue 99554 rewritten
in a-mir-formality, where a single \emph{program} consists of a list of crates
and a type parameter is introduced with the keyword \C{ty}.
%
While rewriting the code, we found that it currently supports only a subset of
Rust and MiniRust;
%
it lacks support for implied bounds, \C{impl} traits, trait objects, function
pointer types, closures, asynchrony, array types, and casts.
%
Furthermore, the \C{?} operator cannot be used because it exists only at the
surface level and is desugared when code is lowered to MIR.
%
As a result, we could express code snippets for only five issues (99554, 98117,
106040, 109628, and 114061).
%
In our experiments, a-mir-formality rejected three (99554, 98117, and 109628)
but accepted the others.
%
Therefore, we conclude that a-mir-formality is also currently immature for
practical use.

To leverage its strong potential as an oracle, we suggest supporting more
features.
%
We believe important features to prioritize include implied bounds, \C{impl}
traits, and trait objects.

\subsection{Documentation}

In this section, we describe \emph{holes} in the existing documentation,
focusing primarily on the Rust Reference~\cite{rust-reference} and the
FLS~\cite{fls}, while also referring to Rust RFCs~\cite{rust-rfcs} where
appropriate.
%
Both the Reference and the FLS describe the syntax in BNF and provide brief
explanations for each syntactic construct.
%
While the Reference has been maintained as an official Rust project for many
years, the FLS was initially developed by Ferrous Systems and AdaCore for use in
safety-critical industries before being adopted as a Rust project in
2025~\cite{adopting-fls}.
%
Although their explanations overlap in many areas, the level of detail provided
for specific features often varies between them.
%
On the other hand, RFCs comprise a collection of proposed changes to the
language.
%
As discussed below, certain features are explained in detail only within RFCs
even after their stabilization.
%
While investigating these documents, we reviewed all sections containing
relevant keywords by utilizing the search functionality on their websites,
avoiding incorrectly claiming that certain information is missing.

Note that the purpose of this discussion is to demonstrate that constructing a
mechanized specification for Rust cannot be achieved simply by collecting
existing documents;
%
rather, it requires significant effort to fill these holes.
%
Our goal is not to criticize existing documents;
%
they were not intended to serve as specifications and remain valuable for
providing intuitive understanding.
%
Since we only aim to demonstrate that the current documentation contains
\emph{some} holes, our discussion does not systematically cover every language
feature.

% Additionally, our goal is not to claim that soundness issues occur due to
% insufficient documentation.
%
% We refer to relevant issues in our discussion to illustrate that the current
% documentation does not provide the precise information required to construct a
% test oracle capable of detecting the collected issues.

{\it Implied bounds and well-formedness.}
%
The Reference explains implied bounds by stating that ``lifetime bounds required
for types to be well-formed are sometimes inferred,'' but it does not define
what well-formed types are.
%
Similarly, the FLS does not provide a precise explanation of implied bounds,
stating only that ``an implied bound is \ldots the byproduct of relations
between lifetime parameters and function parameters \ldots'' without further
clarification on how such relations introduce bounds as a byproduct.
%
The definition of well-formedness exists only in RFCs, notably RFC 1214
\C{projections-lifetimes-and-wf}~\cite{rfc1214}.
%
This RFC defines well-formedness using inference rules based on its own syntax
for Rust types.
%
Unfortunately, this syntax lacks several types present in Rust, most notably
\C{impl} trait types.
%
Furthermore, some inference rules are unclear;
%
e.g., one rule includes the premise \C{('x: 'a) in Env}, yet \C{Env} is never
formally defined.

% Implied bounds serve as a trigger for four issues.
%
% Note that Issue 110925 stems from missing well-formedness checks for \C{impl}
% trait types.
%
% Consequently, while the current definition of well-formedness aids intuitive
% understanding, it remains insufficient for a mechanized specification.

{\it Coercion.}
%
In Issue 99158, coercion from a function item to a function pointer breaks type
soundness.
%
The Reference defines this kind of coercion vaguely, stating only that
%
``coercion is allowed \ldots [from] function item types to \C{fn} pointers.''
%
% although it would be nonsensical to interpret this as allowing the coercion of
% any function item to any function pointer type.
%
The explanation in the FLS is more refined, specifying that coercion is
permitted when
%
``the source's function signature is a subtype of the target's function
signature.''
%
The FLS further defines the subtyping of function pointer types (presumably
including function item signatures) using the standard notion of variance, where
they are contravariant in their parameter types and covariant in their return
types.
%
However, this definition does not cover higher-rank cases, where functions have
type or lifetime parameters, possibly accompanied by bounds.
%
Consequently, it remains unclear whether coercing generic functions to function
pointers requires specific checks for their bounds or permits the substitution
of these parameters with concrete types or lifetimes.

{\it HRTBs and associated types.}
%
In Issue 141713, the use of the type path \C{F::Output} triggers the problem,
where \C{F} has an HRTB.
%
We searched the documentation using the keywords ``higher-rank,'' ``associated
type,'' and ``type path'' to determine how an associated type projected from a
type parameter with an HRTB should be interpreted.
%
Unfortunately, we found explanations only for HRTBs and associated types in
isolation, with no documentation regarding their interaction.
%
Notably, the description of how a type path is resolved to a type is
insufficient even when HRTBs are not involved.
%
The Reference defines the syntax of type paths in its ``Paths'' section but does
not explain their resolution process.
%
The FLS includes an ``Entities and Resolution'' section, which states that \C{<S
as T>::A} ``resolves to the associated type of \C{S}'s trait implementation of
trait \C{T}.''
%
However, it provides no information regarding the resolution of type paths that
lack such explicit trait disambiguation.

%{\it Type inference.}
%%
%The Reference does not precisely describe when type inference succeeds.
%%
%Upon searching for the keyword ``infer,'' we found only two relevant sentences:
%%
%``infer the type if possible based on the surrounding information available''
%and
%%
%``infer the type, or signal an error if insufficient type information is
%available for definite inference.''
%%
%It is not even clear whether there must be only one possible type, although the
%phrase ``definite inference'' may imply this.
%%
%Moreover, it does not explain how ``information'' is collected.

%{\it Arrays.}
%%
%In Issue 144045, the compiler does not check the lifetime of the expression
%\C{[\&x; 0]}.
%%
%The Reference does not specify how expressions of the form \C{[e; N]} are
%type-checked;
%%
%it only states that ``an array is a fixed-size sequence of \C{N} elements of
%type \C{T},'' where ``the array type is written as \C{[T; N]}.''
%%
%This natural-language description is ambiguous.
%%
%We could interpret this to mean that \C{e} must be of type \C{T} for \C{[e; 0]}
%to be typed as \C{[T; 0]} because there is no explicit description for cases
%where \C{N} is zero.
%%
%Alternatively, we might interpret \C{[T; 0]} as a type denoting any empty array,
%%
%although this even allows using the expression \C{[42; 0]} as a value of type
%\C{[\&'static i32; 0]}, which is rejected by the current compiler.

% While five issues are caused by arbitrary type selection during type
% inference,
